\section{Introduction}

In computational numerical linear algebra,
dense matrices often present a computational bottleneck,
scaling quadratically in terms of memory usage ($O(nm)$) for a matrix
with $n$ rows and $m$ columns
and cubically in terms of time complexity ($O(n^3)$) for standard operations.
However, in many real-world applications such as signal
processing\cite{7400965}\cite{signals5030021}
and classical simulation of quantum
systems\cite{10.1098/rsos.160906}, these dense matrices are not totally random.
The elements are arranged in repeating, deterministic patterns.
Designing algorithms that take advantage of this pattern can bypass
the fundamental
limitations of dense linear algebra.

To exploit these patterns, we apply a displacement operation to
determine how much a matrix deviates from a perfect diagonal structure.
For a highly structured matrix, this displacement results in a low rank.
This demonstrates that it is possible to represent a dense matrix
using two smaller, more compact matrices.
Performing operations on these matrices reduces both storage and
computation requirements.
The aim of this project is to explore and implement elementary
operations on quasi-Toeplitz matrices using these displacement operators.
